\section{Approximation and removal of jumps}\label{sec:jumps}

The lower bound in Proposition~\ref{prop:lower} involves a monotone
\(BV\) potential. We first approximate its cost by temperatures with
finitely many jumps. Concavity will then remove those jumps.
For a nonincreasing piecewise \(H^1\) temperature \(u\), equal to 1
inside \(B_1\) and zero after radius \(R\), set
\begin{equation}\label{eq:J}
 J(u)=\int_1^R A(r)|u'(r)|^2\dd r+
       \sum_s A(s)\bigl[\Theta(u(s-))-\Theta(u(s+))\bigr].
\end{equation}
The sum runs over its finitely many downward jumps, including a possible
jump at 1 and the terminal jump to zero. This is an auxiliary cost:
interior jumps are not admissible in the original \(H^1\) problem.

\begin{proposition}\label{prop:approximation}
For every augmented monotone profile \(w\) in \eqref{eq:F} and every
\(\epsilon>0\), there is a temperature of the above form satisfying
\[
 J(u)\le F(w)+\epsilon .
\]
\end{proposition}

\begin{proof}
All approximation parameters are chosen with \(n,\beta,\gamma,R\)
fixed. Write \(A_*=R^{n-1}\), \(L_A=(n-1)R^{n-2}\), and
\(L_G=\operatorname{Lip}(G)\).

\emph{Polygonal potentials.}
Choose a partition of \([1,R]\), of mesh \(h\), with all interior
vertices at continuity points of \(w\). Join the endpoint values
linearly on each cell, using \(w(1+)\) and \(w(R-)\) at the ends,
and keep the two endpoint atoms. Call the result \(w_h\).
For a cell \(C_i=(x_i,x_{i+1})\), let \(d_i\) be its length,
\(y_i\) its left value, and \(m_i\ge0\) its total drop. Thus
\(m_i=\mu(C_i)\), with the retained end atoms excluded, and
\(\sum_i m_i\le Y\).

Freeze \(A_i=A(x_i)\) and \(g_i=G(y_i)\).
For every constant \(p\in[0,1]\),
\[
 pm_i-d_ig_ip^2
   =\int_{C_i}(pz(r)-g_ip^2)\dd r+p\mu_s(C_i)
   \le\int_{C_i}\psi_{g_i}(z(r))\dd r+\mu_s(C_i).
\]
Taking the supremum proves
\[
 d_i\psi_{g_i}(m_i/d_i)
       \le\int_{C_i}\psi_{g_i}(z(r))\dd r+\mu_s(C_i).
\]
For both profiles, the potential stays between \(y_i-m_i\) and \(y_i\).
Unfreezing the weight costs at most \(L_Ahm_i\), since
\(\psi\le z\) and the singular coefficient is 1.
By \eqref{eq:psilip}, unfreezing \(G\) costs at most
\(A_*L_Gd_im_i\). Apply these estimates on both sides and sum:
\begin{equation}\label{eq:polygonalerror}
 F(w_h)\le F(w)+2h(L_A+A_*L_G)Y .
\end{equation}
The endpoint atom costs are identical. This estimate does not require
bounded slopes of \(w\), and applies also to Cantor variation.

\emph{Cutoff near zero.}
Fix \(0<\delta<Y\). Follow the polygonal profile until its potential
first falls to or below \(\delta\), and replace the remainder by zero.
At a continuous crossing this adds a terminal drop \(\delta\).
If an existing jump skips the level, its right value is
\(c\le\delta\); replacing that jump by a drop to zero adds only \(c\).
The same reasoning applies at 1 and \(R\). Removed energy is nonnegative.
For the resulting \(q\),
\begin{equation}\label{eq:cutofferror}
 F(q)\le F(w_h)+A_*\delta.
\end{equation}
If the cutoff is immediate at 1, its cost is the bare-ball cost \(Y\).
Otherwise all positive one-sided values of the retained finite polygon
lie in \([\delta,Y]\), up to its terminal jump at some
\(\ell\in(1,R]\).

\emph{Finite jumps and affine segments.}
Keep the existing jumps of \(q\). Let \(D<\infty\) bound its finitely
many affine slopes, and put
\(g_*=\min_{[\delta,Y]}G>0\), \(g^*=\max_{[\delta,Y]}G\).
Subdivide each affine piece into cells of mesh \(\eta\).
If a cell has left endpoint \(x_i\), length \(d_i\), left value \(y_i\), and drop
\(m_i=z_i d_i\le D\eta\), put
\[
 A_i=A(x_i),\qquad g_i=G(y_i),\qquad b_i=\min(z_i,2g_i).
\]
Insert at the left endpoint a potential drop \(m_i-b_id_i\), then
use an affine segment of slope \(-b_i\) reaching the prescribed
right value. The new potential stays in the cell's original range.
Writing it as \(\widetilde q\), the temperature \(u=Q(\widetilde q)\)
is piecewise \(H^1\), because \(Q'\) is bounded on \([\delta,Y]\).
There are finitely many jumps.

The frozen cell cost is exactly
\begin{equation}\label{eq:microjump}
 A_i\left\{\frac{d_ib_i^2}{4g_i}+m_i-b_id_i\right\}
                   =A_id_i\psi_{g_i}(z_i).
\end{equation}
On its affine part,
\(|u'|^2=b_i^2/(4G(\widetilde q))\),
\(|\widetilde q-y_i|\le m_i\), and \(b_i\le2g^*\).
The error between the actual and frozen Dirichlet terms is at most
\[
 \left(\frac{L_A\eta}{4g_*}
              +\frac{A_*L_Gm_i}{4g_*^2}\right)b_i^2d_i .
\]
The inserted jump has the exact weight \(A_i\).
The difference between the right side of \eqref{eq:microjump}
and the cost of the original affine cell of \(q\) is at most
\(L_A\eta m_i+A_*L_Gd_im_i\).
Since \(m_i\le D\eta\), \(\sum d_i\le R-1\), and
\(\sum m_i\le Y\), the sum of these errors is \(C_\delta\eta\)
for a finite constant depending on the fixed polygon, cutoff and
original parameters. Retained jumps have the same prescribed drop;
if jumps coincide, their potential differences add.
Thus
\[
 J(u)\le F(q)+C_\delta\eta .
\]
Choose first \(h\), then \(\delta\), and finally \(\eta\) to make
the total errors in this inequality, \eqref{eq:polygonalerror}, and
\eqref{eq:cutofferror} less than \(\epsilon\). No uniform bound on
\(1/G\) as \(\delta\downarrow0\) is used.
\end{proof}

\begin{proposition}\label{prop:elimination}
Every finite-jump temperature in \eqref{eq:J} can be replaced, without
increasing its cost, by an admissible temperature for
\((B_1,B_\ell)\) for some \(\ell\in[1,R]\). Its physical energy is
at most \(S J(u)\).
\end{proposition}

\begin{proof}
Take the first downward jump with positive right temperature. Denote
its left and right temperatures by \(b>c>0\), and its radius by \(s\).
On the whole tail after \(s\), replace the potential by
\(\lambda\Theta(u)\), where
\[
 0\le\lambda\le M:=\frac{\Theta(b)}{\Theta(c)}>1.
\]
Its temperature is \(g_\lambda(u)=Q(\lambda\Theta(u))\).
All tail temperatures are at most \(c\), so the inverse is evaluated
in \([0,\Theta(b)]\subset[0,Y]\).

For each fixed \(\lambda>0\), \(g_\lambda\) is Lipschitz on \([0,c]\).
Indeed, if \(\gamma>0\), the leading term of \(\Theta(t)\) at zero
is \(2\beta\gamma^3t^2\); if \(\gamma=0\), it is \(\beta t^5/5\).
Writing the respective order as \(m=2\) or \(m=5\), one obtains
\[
 g_\lambda(t)/t\longrightarrow\lambda^{1/m},\qquad
 g_\lambda'(t)=\frac{\lambda\Theta'(t)}
                         {\Theta'(g_\lambda(t))}
                         \longrightarrow\lambda^{1/m}.
\]
The derivative is bounded away from that endpoint as well.
For \(\lambda=0\) take the zero map. Thus every transformed piece
is \(H^1\), and both endpoint costs are finite.

For \(y=\Theta(u)>0\), the transformed Dirichlet density is
\[
 \frac{\lambda^2|y'|^2}{4G(\lambda y)}
                   =\frac{|y'|^2}{4y^2}H(\lambda y).
\]
At \(\lambda=0\) its value is zero. On \(y=0\), \(y'=0\) a.e.,
and the density is also set to zero there.
Lemma~\ref{lem:concavity} makes this expression concave in \(\lambda\).
Integration preserves that inequality, with finite nonnegative
terms. Every later jump cost, including the terminal one, is linear
in \(\lambda\). At the selected jump the cost is
\[
 A(s)\bigl[\Theta(b)-\lambda\Theta(c)\bigr],
\]
also affine. The total cost \(J_\lambda\) is concave on \([0,M]\);
therefore
\[
 J_1\ge(1-1/M)J_0+J_M/M.
\]
One of the two endpoint choices costs no more than the original profile.

At zero, the entire tail disappears and the selected jump becomes
the outer boundary. At \(M\), the selected jump disappears because
its right temperature becomes \(b\). The number of jumps with
positive right value strictly decreases; monotonicity prevents
new upward jumps. After finitely many repetitions only a terminal
jump to zero remains, at some \(\ell\in[1,R]\).
If the profile reaches zero continuously, this terminal jump has
size zero. Equal traces at the finitely many intervening pieces
cancel their distributional boundary terms. The resulting radial
temperature, extended by 1 in \(B_1\), is in \(H^1(B_\ell)\).
Its gradient and outer trace costs are
\[
 S\int_1^\ell A(r)|u'(r)|^2\dd r
              +S A(\ell)\Theta(u(\ell-))=S J(u).
\]
At \(\ell=1\) this is exactly the bare ball. In particular, the
comparison is now with the original physical energy.
\end{proof}

\begin{proof}[Proof of Theorem~\ref{thm:main}]
Fix any original admissible \(v\). Proposition~\ref{prop:lower}
gives \(w\) with \(S F(w)\le\mathcal E_\Omega(v)\).
For every \(\epsilon>0\), Proposition~\ref{prop:approximation}
gives \(J(u)\le F(w)+\epsilon\), and
Proposition~\ref{prop:elimination} gives an admissible concentric-ball competitor.
Lemma~\ref{lem:radial} ensures that the radial minimum exists.
Hence
\[
 \min_{1\le r\le R}E_\Theta(B_1,B_r)
        \le S J(u)\le S F(w)+S\epsilon
        \le\mathcal E_\Omega(v)+S\epsilon.
\]
Let \(\epsilon\downarrow0\), with the original parameters and
temperature fixed, then minimize over \(v\).
This proves the claimed inequality for the original measurable
cores and bounded Lipschitz outer sets, including the bare ball
and every smaller outer volume.
\end{proof}
